\documentclass{situationist-preprint}

\title{Mary's Room and the Ability Hypothesis}
\msid{MND-04-a}
\msfamily{MND-04}
\msdiscipline{Philosophy of Mind}
\mstype{Principal manuscript}
\mslean{lean/Philosophy/MND/MND\_04.lean}
\mskeywords{knowledge argument; Mary; ability hypothesis; phenomenal concepts; physicalism; qualia}

\begin{document}
\maketitle

\begin{abstract}
Mary knows every physical fact about colour vision but has never seen colours. When she leaves her black-and-white room and sees red, does she learn a new fact? Jackson argued that she does, and that physicalism is therefore false. We argue that she learns no new fact. She acquires new abilities, to recognize, imagine and remember red, and with them new concepts under which she can think of facts she already knew. We show that the knowledge argument is valid only if knowing a fact under a new concept counts as learning a new fact, and that this reading makes the argument's conclusion compatible with physicalism.
\end{abstract}

\section{The argument}

\citet{Jackson1982} describes Mary, a brilliant scientist who has been confined from birth to a black-and-white room and has learned, through black-and-white books and television, everything physical there is to know about colour vision. When she is released and sees a ripe tomato, Jackson argues, she learns something: what it is like to see red. Since she already knew every physical fact, what she learns is not a physical fact, and physicalism is false. The argument can be stated as follows.
\begin{enumerate}[label=(K\arabic*)]
\item Before her release, Mary knows every physical fact about colour vision.
\item On her release, Mary learns a new fact about colour vision.
\item Therefore, there are facts about colour vision that are not physical facts.
\end{enumerate}
The argument is valid, if \enquote{learns a new fact} means that she comes to know a fact she did not know before. Physicalist responses deny (K2), deny (K1), or deny that the argument is valid on the reading on which (K2) is true. Jackson himself later abandoned the argument, coming to think that its conclusion conflicts with the causal efficacy of experience \citep{Jackson1998}.

\subsection{The responses}

The literature on the knowledge argument is large, and the responses can be sorted by the premise they reject \citep{LudlowNagasawaStoljar2004}.

\paragraph{Denying (K1).} \citet{Dennett1991} argued that if Mary really knew everything physical, she would not be surprised: shown a blue banana on her release, she would recognize the trick, because she would know exactly what states a ripe banana would put her brain in and could tell that this was not one of them. On this view the thought experiment trades on our inability to imagine what complete physical knowledge would be like. The response has the merit of taking (K1) seriously, but it is widely felt to beg the question, since whether complete physical knowledge would include knowledge of what experiences are like is what the argument is about.

\paragraph{The ability hypothesis.} \citet{Lewis1988} and \citet{Nemirow1990} held that Mary gains abilities, not information. We examine this response in detail below.

\paragraph{The acquaintance hypothesis.} \citet{Conee1994} argued that Mary gains acquaintance with a property, a kind of knowledge that is neither propositional nor practical. The view shares with the ability hypothesis the claim that Mary learns no new fact, and differs in its account of what she gains.

\paragraph{Old fact, new guise.} \citet{Horgan1984} and others held that Mary learns an old fact in a new way. The phenomenal concept strategy is the most developed version of this response \citep{Loar1990,Papineau2002}: Mary comes to know, under a phenomenal concept, a fact she already knew under a physical one.

\paragraph{Accepting the argument.} Jackson originally drew the dualist conclusion, and combined it with epiphenomenalism about qualia \citep{Jackson1982}. He later abandoned it, concluding that the intuition that Mary learns a new fact must be mistaken, since experiences that made no causal difference could not be known about at all \citep{Jackson1998}.

The view defended here combines the second and fourth responses. The ability hypothesis identifies what Mary gains in practical terms; the phenomenal concept strategy explains why what she gains also has a propositional aspect.

\section{The ability hypothesis}

\citet{Lewis1988} and \citet{Nemirow1990} proposed that what Mary gains is not knowledge that but knowledge how: the abilities to recognize red when she sees it, to imagine it, and to remember it. On this view (K2) is false: Mary learns how to do things, not that something is the case.

\begin{definition}\label{d:ability}
A subject has a \emph{recognitional ability} for a type of experience if she can tell, by having it, that an experience is of that type.
\end{definition}

\begin{definition}\label{d:concept}
A \emph{phenomenal concept} of a type of experience is a concept whose possession consists in, or requires, a recognitional ability for experiences of that type \citep{Loar1990}.
\end{definition}

The ability hypothesis is incomplete as stated. After her release, Mary can think \emph{that experience is like this}, and the thought seems to be a thought that something is the case. The phenomenal concept strategy completes the hypothesis. Mary's new abilities give her new concepts, and with them new ways of thinking about facts she already knew.

\section{The argument refuted}

\begin{assumption}[Physicalism]\label{a:phys}
Every fact about colour vision is a physical fact.
\end{assumption}

\begin{lemma}\label{l:same}
Before and after her release, Mary knows the same facts about colour vision.
\end{lemma}

\begin{proof}
By (K1), before her release Mary knows every physical fact about colour vision. By Assumption~\ref{a:phys}, these are all the facts about colour vision. After her release she knows at least these facts, and there are no others to know.
\end{proof}

\begin{theorem}\label{t:main}
Mary learns no new fact. What she acquires on her release is a recognitional ability and a phenomenal concept, under which she can think of a fact she already knew.
\end{theorem}

\begin{proof}
By Lemma~\ref{l:same}, she knows no new fact. By the description of the case, she acquires the ability to recognize red when she sees it, and by Definition~\ref{d:concept}, a phenomenal concept of seeing red. The thought \emph{seeing red is like this} uses the new concept to think of the physical fact about what happens in her brain when she sees red, which she already knew under a physical concept.
\end{proof}

Theorem~\ref{t:main} does not beg the question against Jackson. It shows that, given physicalism, there is a coherent description of Mary's case on which she undergoes a genuine cognitive change without learning a new fact. Jackson's argument needs (K2) on the reading on which a new fact is learned; physicalists can grant (K2) on the reading on which a new concept is acquired. The argument is valid only on the first reading, and the first reading is what is in dispute.

\section{Variations on the case}

Two variations test the account.

\paragraph{The blue banana.} On her release, Mary's captors show her a blue banana. Does she know that the colour she sees is not yellow? On the present account, she has a new recognitional ability but no phenomenal concept of yellow, since she has never seen yellow. She can tell that the experience she is having is of some colour, and if Dennett is right that her physical knowledge includes knowledge of what her brain does when she sees yellow, she can tell by introspective reports of her brain state that it is not yellow. If Dennett is wrong, she cannot. Either way, what she cannot do is recognize the colour as yellow by its look, and this is what the ability hypothesis predicts.

\paragraph{Marianna.} \citet{NidaRumelin1996} describes Marianna, who is released from the black-and-white room into a room with abstract coloured patches but no coloured objects. She acquires phenomenal concepts of the colours, but does not know which colour words they correspond to: she cannot tell whether the patch she is looking at is the colour that the sky normally has. When she later sees the sky, she learns something. On the present account, what she learns is an identity between a phenomenal concept and a physical one: that the experience she has been calling \emph{this} is the experience produced by light of a certain wavelength. Like the discovery that Hesperus is Phosphorus, this is genuine learning, but it is learning of an identity between two ways of thinking of one fact, not of a new fact. Marianna's case is therefore the clearest illustration of the old fact, new guise response.

\section{Objections to the ability hypothesis}

\begin{objection}
Mary's new knowledge embeds in propositional contexts. She can come to believe that if seeing red is like \emph{this}, then seeing pink is like \emph{that}, and can reason from such beliefs. Abilities do not embed in conditionals \citep{Loar1990}.
\end{objection}

\begin{reply}
The objection tells against the ability hypothesis in its pure form, on which Mary gains only know-how. It does not tell against the combined view defended here, on which her abilities give her phenomenal concepts, and the concepts can occur in propositional thoughts, including conditional ones. The conditional thought uses phenomenal concepts to think about physical facts she already knew.
\end{reply}

\begin{objection}
Abilities are neither necessary nor sufficient for knowing what an experience is like. Someone with a poor visual imagination may know what red looks like while being unable to imagine it, and someone could acquire the abilities by brain surgery without having had the experience \citep{Alter1998}.
\end{objection}

\begin{reply}
The ability that matters is recognitional, not imaginative, which answers the first case: the person with a poor imagination can still recognize red when she sees it. The second case is one in which the subject acquires a phenomenal concept by an unusual route. On the present account she then knows what seeing red is like, as she should: what matters for possessing the concept is the recognitional capacity, however acquired.
\end{reply}

\section{An empirical analogue}

Mary's case is a thought experiment, but it has an empirical relative in Molyneux's problem: would a person born blind, who had learned to distinguish a cube from a sphere by touch, be able to tell them apart by sight alone on gaining vision? Locke and Molyneux thought not. \citet{Held2011} tested the question with patients who gained sight after surgery for congenital cataracts. Immediately after surgery, the patients could not match objects seen to objects felt at better than chance, although they could match seen objects with seen objects; within days they acquired the ability.

The finding bears on Mary's case in a limited but suggestive way. It shows that recognitional abilities in a new modality are not derivable from knowledge acquired in another, and must be acquired through experience in the new modality. This is what the ability hypothesis predicts for Mary: her complete physical knowledge, acquired through black-and-white media, does not give her the recognitional ability for colour, which she must acquire by seeing. The finding is neutral on whether what is acquired is a new fact or a new ability, and so does not decide the debate; but it removes any suggestion that Mary's inability before her release is a peculiarity of the thought experiment.

\section{Relations to other results}

The view defended here is the one that underlies several other results in the collection. MND-03 argues that much of what it is like to be a bat can be specified in structural terms, but not all of it; the residue is what can be grasped only through phenomenal concepts, which humans lack for echolocation as Mary lacks them for colour. MND-10 argues that qualia exist and are intrinsic, although their structure can be described relationally; on the present view, phenomenal concepts are the concepts under which qualia are grasped directly, and physical concepts the concepts under which their structure is described. MND-05 applies the same distinction to the zombie argument.

\section{Formalization}

The module \nolinkurl{Philosophy/MND/MND_04.lean} represents facts as objects, with a predicate for being physical, and knowledge before and after release as predicates on facts. It proves Lemma~\ref{l:same} from (K1) and Assumption~\ref{a:phys}, and shows that the conclusion of the knowledge argument is inconsistent with physicalism only given the reading of (K2) on which a new fact is known.

\begin{lstlisting}
theorem no_new_fact (f : Fact) : KnowsAfter f → KnowsBefore f
\end{lstlisting}

\section{Objections}

\begin{objection}
Mary is surprised. If she learned nothing new, why would she be surprised?
\end{objection}

\begin{reply}
Surprise attaches to modes of presentation as well as to facts. One can be surprised to learn that Hesperus is Phosphorus, although one has learned no new fact about the planet. Mary is surprised because she encounters, for the first time, a phenomenal mode of presentation of a fact she knew under a physical one.
\end{reply}

\begin{objection}
If Mary already knew everything physical, she should have been able to work out what seeing red is like. Her inability shows that this knowledge is not physical.
\end{objection}

\begin{reply}
Her inability shows that physical knowledge does not supply phenomenal concepts, which require the corresponding experience. It does not show that the facts the concepts pick out are non-physical. \citet{Dennett1991} argued that Mary could work it out; we need not go so far. The relation between physical and phenomenal concepts is a posteriori, like the relation between \enquote{water} and \enquote{H$_2$O}, and a posteriori identities cannot be derived a priori from one side.
\end{reply}

\section{Conclusion}

Mary learns no new fact. She acquires abilities and, with them, concepts under which she thinks of facts she already knew. The knowledge argument trades on an ambiguity in \enquote{learning something new}, and on the reading on which its premise is true, its conclusion does not follow. The companion results on zombies (MND-05) and on the hard problem (MND-01) rest on the same distinction between concepts and facts.

\begin{thebibliography}{99}\small
\bibitem[Alter(1998)]{Alter1998} Alter, T. (1998). A limited defense of the knowledge argument. \emph{Philosophical Studies}, 90(1), 35--56.
\bibitem[Conee(1994)]{Conee1994} Conee, E. (1994). Phenomenal knowledge. \emph{Australasian Journal of Philosophy}, 72(2), 136--150.
\bibitem[Dennett(1991)]{Dennett1991} Dennett, D.~C. (1991). \emph{Consciousness Explained}. Little, Brown.
\bibitem[Held et~al.(2011)]{Held2011} Held, R., Ostrovsky, Y., de~Gelder, B., Gandhi, T., Ganesh, S., Mathur, U. and Sinha, P. (2011). The newly sighted fail to match seen with felt. \emph{Nature Neuroscience}, 14(5), 551--553.
\bibitem[Horgan(1984)]{Horgan1984} Horgan, T. (1984). Jackson on physical information and qualia. \emph{The Philosophical Quarterly}, 34(135), 147--152.
\bibitem[Jackson(1982)]{Jackson1982} Jackson, F. (1982). Epiphenomenal qualia. \emph{The Philosophical Quarterly}, 32(127), 127--136.
\bibitem[Jackson(1998)]{Jackson1998} Jackson, F. (1998). Postscript on qualia. In \emph{Mind, Method, and Conditionals: Selected Essays} (pp.~76--79). Routledge.
\bibitem[Lewis(1988)]{Lewis1988} Lewis, D. (1988). What experience teaches. \emph{Proceedings of the Russellian Society}, 13, 29--57.
\bibitem[Loar(1990)]{Loar1990} Loar, B. (1990). Phenomenal states. \emph{Philosophical Perspectives}, 4, 81--108.
\bibitem[Ludlow et~al.(2004)]{LudlowNagasawaStoljar2004} Ludlow, P., Nagasawa, Y. and Stoljar, D. (Eds.) (2004). \emph{There's Something About Mary: Essays on Phenomenal Consciousness and Frank Jackson's Knowledge Argument}. MIT Press.
\bibitem[Nemirow(1990)]{Nemirow1990} Nemirow, L. (1990). Physicalism and the cognitive role of acquaintance. In W.~G. Lycan (Ed.), \emph{Mind and Cognition: A Reader} (pp.~490--499). Blackwell.
\bibitem[Nida-Rümelin(1996)]{NidaRumelin1996} Nida-Rümelin, M. (1996). What Mary couldn't know: Belief about phenomenal states. In T. Metzinger (Ed.), \emph{Conscious Experience} (pp.~219--241). Imprint Academic.
\bibitem[Papineau(2002)]{Papineau2002} Papineau, D. (2002). \emph{Thinking about Consciousness}. Oxford University Press.
\end{thebibliography}

\end{document}
